\chapter{Redes de Computadores}
\label{cap:redes}

\begin{pegadinha}[Fora do edital do Perfil 3, mas cai --- não pule este capítulo]
Redes de Computadores \textbf{não está} no conteúdo do Perfil 3 (está nos
Perfis 2 e 5). O edital do Perfil 3 só cita HTTPS e SSL/TLS. Mesmo assim, a
\textbf{Dataprev 2024 trouxe ${\sim}$3 questões de redes} para Desenvolvimento
de Software (modelo OSI, protocolo, segurança de rede). Ignorar isso custou
7,5 pontos. Estude o essencial abaixo --- é ponto barato e seguro.
\end{pegadinha}

\section{Modelos de referência}

\textbf{OSI (7 camadas)} --- de baixo para cima:

\begin{center}
\begin{tabular}{@{}p{0.6cm}p{2.2cm}p{4cm}p{3cm}@{}}
\toprule
\# & \textbf{Camada} & \textbf{Função} & \textbf{Exemplos} \\
\midrule
1 & Física & bits no meio & cabos, sinais \\
2 & Enlace & quadros, MAC & switch, Ethernet \\
3 & Rede & pacotes, IP, roteamento & roteador, IP \\
4 & Transporte & fim a fim, portas & TCP, UDP \\
5 & Sessão & diálogo, checkpoints & --- \\
6 & Apresentação & formato, cifra, compressão & TLS (discutível) \\
7 & Aplicação & serviços ao usuário & HTTP, DNS, SMTP \\
\bottomrule
\end{tabular}
\end{center}

\textbf{TCP/IP (4 camadas):} Acesso à Rede, Internet (IP), Transporte
(TCP/UDP), Aplicação.

\begin{pegadinha}
Modelo OSI é cobrado por \textbf{função}, não por número: o cenário descreve
o comportamento e pede a camada. Ex.: ``checkpoints no fluxo + controle de
diálogo (simplex/half/full-duplex)'' = camada de \textbf{Sessão} --- não
Transporte, Enlace ou Apresentação, que são os distratores comuns. Troca
clássica de elemento por camada: switch = 2 (enlace), roteador = 3 (rede).
\end{pegadinha}

\section{Equipamentos}

\textbf{Hub} (camada 1, burro, repete para todos) $\times$ \textbf{Switch}
(camada 2, usa MAC, comuta para a porta certa) $\times$ \textbf{Roteador}
(camada 3, usa IP, interliga redes). Existe switch L3 (roteia), mas o par
clássico cobrado é switch=2, roteador=3.

\section{TCP $\times$ UDP}

\begin{center}
\begin{tabular}{@{}p{2.6cm}p{4.6cm}p{4.6cm}@{}}
\toprule
& \textbf{TCP} & \textbf{UDP} \\
\midrule
Conexão & orientado a conexão (handshake 3 vias) & sem conexão \\
Confiabilidade & garante entrega e ordem & não garante \\
Uso & web, e-mail, transferência & streaming, DNS, VoIP \\
\bottomrule
\end{tabular}
\end{center}

\begin{pegadinha}
O \textbf{three-way handshake} do TCP \textbf{sincroniza números de
sequência} (SYN / SYN-ACK / ACK). Os distratores dizem que ele ``autentica
dispositivos'', ``negocia segurança'' ou ``controla fluxo'' --- nenhum dos
três é o papel do handshake. Handshake existe \textbf{só no TCP}.
\end{pegadinha}

\section{Protocolos e portas}

\begin{center}
\begin{tabular}{@{}p{2.6cm}p{2.4cm}p{6.2cm}@{}}
\toprule
\textbf{Protocolo} & \textbf{Porta} & \textbf{Função} \\
\midrule
HTTP & 80 & web \\
HTTPS & 443 & web seguro (TLS) \\
FTP & 21 (20 dados) & transferência \\
SSH & 22 & acesso remoto seguro \\
SMTP & 25 & envio de e-mail \\
DNS & 53 & nomes $\to$ IP \\
POP3 / IMAP & 110 / 143 & leitura de e-mail \\
DHCP & 67/68 & IP automático \\
\bottomrule
\end{tabular}
\end{center}

\textit{Observação:} porta específica não apareceu na amostra de provas
reais analisadas; a FGV preferiu raciocínio de rota, camada e ACL. Está no
edital, mas priorize os padrões desta seção.

\section{Endereçamento e roteamento IP}

\textbf{IPv4:} 32 bits. Faixas privadas (RFC 1918): 10.0.0.0/8,
172.16.0.0/12, 192.168.0.0/16. \textbf{IPv6:} 128 bits; prefixo de sub-rede
típico /64.

\textbf{Roteamento:} rota estática $\times$ dinâmica (OSPF, RIP, BGP).
\textbf{Distância administrativa} (Cisco) desempata a fonte da rota:
conectada 0, estática 1, OSPF 110, RIP 120 --- \textbf{menor} vence. Com duas
rotas default, ganha a estática (1) sobre OSPF (110).

\begin{pegadinha}
Absolutos em roteamento denunciam o distrator: ``OSPF SEMPRE tem prioridade
sobre rota estática'', ``roteador faz load balancing AUTOMÁTICO'', ``descarta
por rotas conflitantes''. A regra real é determinística: menor distância
administrativa vence, sempre.
\end{pegadinha}

\subsection{Sub-redes: máscara $\times$ hosts utilizáveis}

A conta é sempre a mesma: com $h$ bits sobrando para host,
\textbf{hosts utilizáveis} $= 2^h - 2$ (descontam-se o endereço de
\textbf{rede} e o de \textbf{broadcast}). O prefixo $/n$ deixa $h = 32 - n$.

\begin{center}
\begin{tabular}{@{}p{1.6cm}p{3.6cm}p{2.2cm}p{3cm}@{}}
\toprule
\textbf{Prefixo} & \textbf{Máscara} & \textbf{Bits de host} & \textbf{Hosts utilizáveis} \\
\midrule
/24 & 255.255.255.0 & 8 & 254 \\
/25 & 255.255.255.128 & 7 & 126 \\
/26 & 255.255.255.192 & 6 & 62 \\
/27 & 255.255.255.224 & 5 & \textbf{30} \\
/28 & 255.255.255.240 & 4 & 14 \\
\bottomrule
\end{tabular}
\end{center}

O enunciado costuma dar a necessidade (``no máximo 30 hosts'') e pedir a
máscara: procure a \textbf{menor} sub-rede que \textbf{comporta} o número
pedido --- 30 hosts cabem em /27 exatamente, então /26 (62) desperdiça e /28
(14) não cabe.

\textbf{IPv6 --- notação simplificada.} Duas regras, nesta ordem: (1) apague
os \textbf{zeros à esquerda} de cada grupo (\texttt{0DB8} $\to$
\texttt{DB8}, \texttt{0001} $\to$ \texttt{1}); (2) substitua \textbf{uma
única} sequência de grupos inteiramente nulos por \texttt{::} --- só uma vez
no endereço, senão fica ambíguo. Assim
\texttt{2001:0DB8:0000:0000:0000:0000:FE00:0001} vira
\texttt{2001:DB8::FE00:1}.

\subsection{ACL Cisco e wildcard mask}

Wildcard mask = inverso da máscara; bit \textbf{0} = ``tem que bater'', bit
\textbf{1} = ``ignora''. Para casar só endereços ímpares, o bit menos
significativo do octeto precisa estar fixo em 1 $\to$ wildcard
\texttt{0.0.0.254} (deixa livres os demais bits, fixa o último). Confira
montando os bits, não de cor.

\subsection{MPLS, VLAN, NAT}

\begin{itemize}
\item \textbf{MPLS (Multiprotocol Label Switching):} encaminha por
\textbf{rótulos (labels)} em vez de olhar o IP de destino a cada salto ---
por isso é chamado de \textbf{``camada 2.5''} (entre enlace/L2 e rede/L3).
Cria caminhos (LSP) para engenharia de tráfego e QoS.
\item \textbf{VLAN:} segmenta logicamente a camada 2 (broadcast domains
separados no mesmo switch físico).
\item \textbf{NAT:} traduz endereços privados $\leftrightarrow$ público
(economiza IPv4, oculta a rede interna).
\end{itemize}

\begin{pegadinha}
MPLS é \textbf{rótulo} (``2.5''), não roteamento IP puro; VLAN segmenta
\textbf{L2}. Não confundir os três termos entre si.
\end{pegadinha}

\section{Tipos de rede corporativa}

Internet (pública global) $\times$ intranet (interna) $\times$ extranet
(acesso controlado a parceiros/clientes externos) $\times$ portal (ponto
único de acesso) --- ver também Capítulo \ref{cap:arquitetura}.

\begin{pegadinha}
Inverter intranet/extranet/Internet: chamar intranet de ``rede pública
global'', extranet de ``só interna'', Internet de ``rede restrita'' --- os
três papéis trocados são o distrator clássico.
\end{pegadinha}

\section{Segurança de rede (interface com Segurança da Informação)}

\textbf{Firewall} (filtra tráfego), \textbf{proxy}, \textbf{VPN} (IPsec, SSL
VPN). \textbf{IDS} (detecta/alerta) $\times$ \textbf{IPS} (bloqueia) --- ver
Capítulo \ref{cap:seguranca}. \textbf{SSL $\times$ TLS $\times$ HTTPS} e
\textbf{X.800} (mecanismos de segurança OSI) também são fronteira com
segurança.

\subsection{Firewall stateful $\times$ stateless}

\begin{center}
\begin{tabular}{@{}p{2.4cm}p{4.6cm}p{4.6cm}@{}}
\toprule
& \textbf{Stateless} (filtro de pacotes) & \textbf{Stateful} (com estado) \\
\midrule
O que olha & cada pacote \textbf{isoladamente}: IP de origem/destino, porta,
protocolo & o pacote \textbf{no contexto da sessão} \\
Como decide & só pela regra estática & mantém uma \textbf{tabela de estado}
das conexões ativas \\
Efeito prático & para liberar a resposta é preciso escrever a regra de volta
à mão & a resposta de uma conexão iniciada de dentro já é aceita
automaticamente \\
\bottomrule
\end{tabular}
\end{center}

Os dois atuam nas camadas 3 e 4. Quem sobe para a camada 7 (inspeciona o
conteúdo da aplicação) é o \textbf{firewall de próxima geração (NGFW)} ou o
\textbf{WAF}, que são outra coisa.

\subsection{SSH $\times$ Telnet e o handshake do TLS}

\textbf{Telnet} (porta \textbf{23}) trafega \textbf{tudo em texto claro},
inclusive a senha do administrador --- quem estiver no caminho lê a sessão
inteira. O \textbf{SSH} (porta \textbf{22}) faz o mesmo trabalho de
administração remota \textbf{cifrando toda a comunicação, inclusive a
autenticação}. É o exemplo canônico de protocolo inseguro aposentado por um
seguro, e a razão da substituição é \textbf{criptografia}, não velocidade.

No \textbf{HTTPS}, o TLS usa criptografia \textbf{híbrida}, e a divisão de
trabalho é o que a banca cobra:

\begin{itemize}
\item \textbf{Assimétrica --- só no handshake.} O certificado digital do
servidor autentica quem ele diz ser, e o par de chaves serve para as duas
pontas \textbf{acordarem com segurança uma chave simétrica de sessão}.
\item \textbf{Simétrica --- no resto.} Todo o volume de dados da sessão é
cifrado com essa chave simétrica, que é \textbf{muito mais rápida}.
\end{itemize}

A assimétrica é lenta demais para cifrar o tráfego inteiro; a simétrica não
resolveria sozinha o problema de \emph{como} combinar a chave por um canal
inseguro. Daí o híbrido.

\begin{pegadinha}
``SSL mais seguro/moderno que TLS'' está \textbf{invertido} --- TLS sucedeu e
corrigiu o SSL.

No handshake TLS, o distrator diz que a assimétrica ``cifra todo o tráfego da
sessão'' (não: ela só \textbf{troca a chave}), que ela ``autentica o usuário
final por login e senha'' (não: autentica o \textbf{servidor}, por
certificado) ou que ``dispensa os certificados'' (é justamente o contrário).

Em firewall, o distrator descreve o \textbf{stateless} e chama de stateful
(``analisa apenas os cabeçalhos de cada pacote isoladamente'') --- inversão
de par pura. Outro: dizer que o stateful ``atua exclusivamente na camada de
aplicação'' --- absoluto e errado, isso é NGFW/WAF.

Em SSH, o distrator inverte o motivo: ``é mais rápido por não usar
criptografia'' ou ``transmite as credenciais em texto claro'' --- essa é a
descrição do \textbf{Telnet}.
\end{pegadinha}

\begin{jacaiu}
\textbf{Em prova real da FGV:} camada OSI \textbf{por função} --- sessão e
\emph{checkpoints} na retomada de uma transferência interrompida ---,
\textbf{rota estática $\times$ OSPF} pela distância administrativa e
\textbf{ACL Cisco com wildcard mask}; ainda volume anômalo de handshakes TCP
(SYN flood) --- \textbf{TJ-RJ}. Camada OSI para \textbf{diagnosticar} falha
entre servidor e switch; \textbf{switch $\times$ roteador $\times$
hub/repetidor}; \textbf{TCP $\times$ UDP} (entrega ordenada e garantida);
\textbf{RFC 1918 e NAT}; \textbf{IPv6} (abreviação do endereço);
\textbf{cálculo de sub-rede} a partir de um /24; \textbf{MPLS}; meio físico,
topologia e Wi-Fi; \textbf{SSH substituindo o Telnet} ---
\textbf{ALERO 2026}. \textbf{Internet $\times$ intranet $\times$ extranet
$\times$ portal} e \textbf{X.800} (mecanismos de segurança do modelo OSI) ---
\textbf{Dataprev 2024}.

\textbf{No nosso banco} (previsto pelo edital, ainda não visto na amostra de
provas): \textbf{número de porta específico} (80/443, DNS 53, DHCP) --- veja
a observação da Seção de protocolos, a FGV preferiu raciocínio de rota,
camada e ACL; o \textbf{handshake híbrido do TLS}; \textbf{firewall stateful
$\times$ stateless}; e \textbf{VLAN} como segmentação de domínio de
broadcast.

Rode \texttt{./quiz.py redes}.
\end{jacaiu}

\begin{comosair}
Não vale estudar redes com a profundidade dos Perfis 2/5. Cubra:
OSI/TCP-IP, equipamentos por camada, TCP$\times$UDP, portas comuns, IP
público$\times$privado, IDS$\times$IPS. Isso resolve o padrão histórico da
FGV para o perfil de desenvolvimento com pouco tempo investido. OSI de cima
pra baixo, para decorar rápido: Aplicação, Apresentação, Sessão, Transporte,
Rede, Enlace, Física.
\end{comosair}
